\chapter{Bloch Frames, Alignment Families, and Fourier Locality}
\label{ch:bloch-alignment-locality}

\section{The comparison problem in finite band reduction}

\subsection{From a parent Hamiltonian to a retained model}

Many electronic-structure and reduced-order calculations begin with an operator that is
too large, too detailed, or expressed in the wrong coordinates for the intended
analysis. On a finite momentum mesh, let $\mathcal K$ denote the supplied set of crystal
momenta and let $N_k=|\mathcal K|$. At each $k\in\mathcal K$, a parent calculation
provides a finite Hermitian matrix
\[
  H_{\mathrm{parent}}(k)\in\mathbb C^{D\times D}.
\]
The dimension $D$ may count plane waves, localized orbitals, finite-difference degrees
of freedom, or any other finite ambient basis; the same algebra applies to each case.

Suppose the scientific question concerns only $M$ states, with $M\leq D$, such as an
isolated band, a composite group of bands near the Fermi level, or a trial subspace used
for downfolding. These states are assembled as the columns of a frame
\[
  V(k)\in\mathbb C^{D\times M}.
\]
When $V(k)^\dagger V(k)=I_M$, the columns are orthonormal and determine the retained
subspace
\[
  \mathcal S(k)=\operatorname{Ran}V(k)
\]
and its ambient orthogonal projector
\[
  P(k)=V(k)V(k)^\dagger.
\]
The parent operator represented in retained coordinates is
\[
  h(k)=V(k)^\dagger H_{\mathrm{parent}}(k)V(k)
  \in\mathbb C^{M\times M}.
\]
This matrix is the finite projected Hamiltonian. Spectral isolation and accuracy for a
target observable are additional properties of the selected subspace.

In applications, the retained frame $V(k)$ may be obtained by diagonalizing the parent
Hamiltonian, projecting trial orbitals, disentangling a composite band manifold,
constructing Wannier functions, or fitting a reduced model. The finite-frame theorem
sequence in Chapter~\ref{ch:finite-operator-foundations} takes the resulting matrix
$V(k)$ and explicit hypotheses such as $V(k)^\dagger V(k)=I_M$ as its inputs.

\subsection{Why two reduced matrices cannot be compared entry by entry}

The columns of $V(k)$ are coordinates, not uniquely labeled physical objects. If
$G(k)\in U(M)$ is unitary, then
\[
  V(k)\longmapsto V(k)G(k)
\]
changes the retained basis but not the subspace $\mathcal S(k)$. The projector remains
$P(k)$, while the compressed Hamiltonian becomes
\[
  h(k)\longmapsto G(k)^\dagger h(k)G(k).
\]
Thus the subspace is invariant, but the retained matrix is only covariant. Its entries
and eigenvectors depend on the chosen retained coordinates even though its spectrum is
unchanged.

This distinction becomes unavoidable when comparing a reference model with a candidate
model. Let
\[
  h_{\mathrm{ref}}(k),\qquad h_{\mathrm{cand}}(k)
  \in\mathbb C^{M\times M}
\]
denote their retained operators. A direct difference
$h_{\mathrm{ref}}(k)-h_{\mathrm{cand}}(k)$ is meaningful only if the two matrices are
already expressed in the same retained coordinates. More generally, an identification
$U(k)$ from candidate coordinates to reference coordinates must be supplied, and the
comparison uses the pulled-back candidate operator
\[
  U(k)^{-1}h_{\mathrm{cand}}(k)U(k).
\]
The sign convention used in this repository is
\[
  \Delta_U(k)
  =h_{\mathrm{ref}}(k)
   -U(k)^{-1}h_{\mathrm{cand}}(k)U(k).
\]
For a unitary identification, $U(k)^{-1}=U(k)^\dagger$. The norm of $\Delta_U(k)$
measures an operator discrepancy relative to that supplied identification. Its
interpretation therefore includes the construction and admissibility of $U$.

Band energies are insufficient for this comparison. Pointwise unitary conjugation preserves
the eigenvalues of $h(k)$, so two fields can have identical band energies while their
matrix-valued eigenvectors, Fourier coefficients, and real-space representations differ.
An adequacy study that cares about operator structure must therefore specify more than a
spectral tolerance.

\subsection{Why momentum dependence affects locality}

A real-space reduced model is obtained by Fourier transforming the momentum-dependent
operator field. In conventional finite-mesh notation, its block at translation $R$ is
\[
  T_R[h]
  =\frac{1}{N_k}\sum_{k\in\mathcal K}
    e^{-\mathrm{i}k\cdot R}h(k).
\]
Here $R$ labels a real-space lattice translation; in the exact Lean development the
phase factors are finite-group characters. Short-range or sparse effective models keep
only a declared set of $R$ values. Consequently, Fourier support and the distribution of
Frobenius power among real-space blocks are part of the model-reduction question, not
merely presentation choices.

If one constant unitary $G$ is used at every $k$, it factors through the finite Fourier
sum:
\[
  T_R[G^\dagger hG]=G^\dagger T_R[h]G.
\]
The real-space labels are not mixed. Exact support truncation therefore commutes with
this common change of retained coordinates, and unitary-invariant block norms are
preserved.

The situation changes when $G$ depends on $k$:
\[
  T_R[G(\cdot)^\dagger h(\cdot)G(\cdot)]
  =\frac{1}{N_k}\sum_{k\in\mathcal K}
    e^{-\mathrm{i}k\cdot R}G(k)^\dagger h(k)G(k).
\]
The matrices $G(k)$ can no longer be factored outside the sum. A pointwise alignment may
therefore improve an operator comparison at each momentum while changing how the same
field is distributed over real-space blocks. This is the tension at the center of the
paper: pointwise spectral or operator agreement and real-space locality are not the same
criterion. Section~\ref{sec:gauge-truncation-counterexample} gives the checked exact
two-point witness: a varying rank-two gauge creates a nonzero odd mode and does not
commute with zero-mode support truncation.

\subsection{The admissible alignment family is part of the claim}

The phrase ``optimize over the gauge'' is incomplete until the allowed gauges are
specified. Several mathematically different questions arise:

\begin{itemize}
  \item Is no coordinate adjustment allowed, so that $U(k)=I_M$?
  \item May one global unitary be used at every momentum?
  \item May an arbitrary unitary be chosen independently at every momentum?
  \item Must $U(k)$ satisfy continuity, endpoint sewing, symmetry, orientation, or a
        finite-depth structural ansatz?
\end{itemize}

These choices define different feasible sets. Enlarging the alignment family can only
make an existential adequacy test easier to satisfy, while the reverse implication
fails in general. A pointwise witness belongs to the pointwise family; transfer to a
global, smooth, or structured family requires an explicit membership proof. The strict
family-separation theorems later in this chapter make this distinction exact.

For that reason, the operative adequacy question is joint and family-relative: does one
candidate, one support choice, and one alignment from the declared family satisfy all
operator, spectral, and locality thresholds? Such a witness uses the same alignment for
every objective. Envelopes, plots, and optimizer traces summarize searches rather than
supply the existential witness or an infeasibility certificate.

\subsection{Role of the finite-frame theorems}

The ten finite-frame results proved in
Chapter~\ref{ch:finite-operator-foundations} supply the algebraic foundation for this
comparison. They show exactly which objects are invariant when only retained coordinates
change, which objects transform covariantly, and how a compressed operator can be
reconstructed as an ambient projector sandwich. This prevents three common category
errors:

\begin{enumerate}
  \item treating a basis change as a change of retained subspace;
  \item comparing compressed matrices before identifying their retained coordinates;
        and
  \item inferring physical or spectral adequacy from algebraic covariance alone.
\end{enumerate}

The present statements are finite and exact. Construction of $V(k)$ and $U(k)$,
continuity, localization, and material interpretation enter through additional
mathematical or evidential layers.

\section{Relation to the research manuscript's Bloch and operator-comparison chapters}

The working composite-band chapter denotes its momentum-dependent frame by
\[
  F(\kappa)=\begin{bmatrix}u_1(\kappa)&\cdots&u_M(\kappa)\end{bmatrix}.
\]
At corresponding momentum coordinates, the notation of this paper is
\[
  V(k)=F(\kappa)
  =\sum_{a=1}^M |u_a(\kappa)\rangle\,{}_{\mathrm{ret}}\!\langle a|.
\]
Thus $V(k)$ and $F(\kappa)$ are two names for the map that inserts retained-coordinate
vectors into an ambient Bloch fiber.

\begin{table}[ht]
\centering
\small
\begin{tabularx}{\textwidth}{@{}l X X@{}}
\toprule
This paper & Research manuscript & Common mathematical object \\
\midrule
$V(k)$
  & $F(\kappa)$
  & momentum-dependent retained-frame embedding \\
$|v_a(k)\rangle$
  & $u_a(\kappa)$
  & column $a$ of the composite Bloch frame \\
$P_V(k)$
  & $P(\kappa)=F(\kappa)F(\kappa)^\dagger$
  & ambient projector field \\
$C_{V(k)}(H(k))$
  & $H_F(\kappa)$
  & retained operator field \\
$T_R[h]$
  & $H_U(R)$ or $\langle w_{m0}|\widehat H|w_{nR}\rangle$
  & real-space Fourier block of the retained operator \\
\bottomrule
\end{tabularx}
\caption{Bloch and Fourier notation bridge to the working research manuscript.}
\label{tab:bloch-manuscript-notation-bridge}
\end{table}

This paper uses $G(k)$ for an internal unitary change of a single retained frame and
reserves $U(k)$ for an identification between reference and candidate retained spaces.
The working composite-band chapter writes its internal gauge as $U(\kappa)$. These
uses must not be merged silently: an internal gauge and a two-model identification have
different domains, provenance, and adequacy obligations.

The aligned-difference sign also differs. When the research manuscript's pristine
operator is identified with $h_{\mathrm{ref}}$, its doped operator with
$h_{\mathrm{cand}}$, and its map $U$ runs from reference to candidate coordinates, it
writes
\[
  \Delta H_{d\to b}=U^\dagger h_{\mathrm{cand}}U-h_{\mathrm{ref}}.
\]
The Formal convention used here is
\[
  \Delta_U=h_{\mathrm{ref}}-U^{-1}h_{\mathrm{cand}}U
  =-\Delta H_{d\to b}
\]
for unitary $U$. Norms are insensitive to this sign, but a signed perturbation or
impurity interpretation is not. Any later editorial export must state the operand order.

This bridge records notation observed at the same clean research-manuscript commit,
principally in \modulepath{chapters/composite-band-gauge-alignment.tex} and
\modulepath{chapters/operator-comparison.tex}. Formal source fixes theorem status.

\section{Aligned operator differences}

The Bloch layer represents a pointwise unitary alignment by an isometric linear
equivalence $U(k)$ between reference and candidate fibers. The candidate operator is
pulled back into the reference fiber before subtraction. Declarations
\decl{BG-BLOCH-THM-012--020} prove pointwise evaluation, closure of shared sewing under
subtraction, sign reversal, recovery, an exact-zero criterion, and sewing closure of the
aligned difference.

The sewing theorem assumes sewing laws for both operator fields and an alignment that
intertwines them. Under these hypotheses, the aligned difference inherits sewing.
Self-adjointness and regularity require their own hypotheses.

\section{Alignment families are decision spaces}

Declarations \decl{BG-ALIGNMENT-DEF-007--016} distinguish:

\begin{itemize}
  \item one declared singleton alignment;
  \item identity-only alignment;
  \item one global unitary reused at every point; and
  \item the complete family of arbitrary pointwise unitary alignments.
\end{itemize}

Feasibility means existence of an accepted alignment inside one specified family.
Declarations \decl{BG-ALIGNMENT-THM-011--019} prove the forward inclusions
\[
  \mathcal U_{\mathrm{id}}
  \subseteq
  \mathcal U_{\mathrm{global}}
  \subseteq
  \mathcal U_{\mathrm{pointwise}}
\]
and transfer feasibility only in that direction.

Declarations \decl{BG-ALIGNMENT-THM-020--029} provide strictness witnesses when the
relevant point and unitary inequalities are supplied. They also construct acceptance
predicates feasible in a larger family but infeasible in a proper subfamily. Therefore a
pointwise-alignment witness need not belong to the global-unitary or identity family.
The result is a logical separation of decision spaces.

\section{Fixed and pointwise gauge actions}

A single unitary common to all finite coordinates preserves Frobenius inner products,
powers, and differences and commutes with finite Fourier coefficient extraction,
reconstruction, exact support truncation, and the corresponding residual. These results
are \decl{BG-FOURIER-THM-033--047}.

A pointwise unitary field has a different boundary. Declarations
\decl{BG-FOURIER-THM-048--054} prove an exact character-phase mode shift, the triviality
of scalar operator conjugation, coefficient separation, reduction of a constant field
to the fixed-unitary case, conditional coefficient change for a genuinely changed
field, and preservation of weighted pointwise Frobenius power.

The last result is invariance of weighted pointwise Frobenius power. Fourier
coefficients, support, decay, and locality require separate statements; a general
pointwise-gauge convolution theorem is outside the current family.

\section{An exact gauge--truncation counterexample}
\label{sec:gauge-truncation-counterexample}

Let the momentum group be $G=\mathbb Z/2\mathbb Z$, and define
\[
  P_0=\begin{bmatrix}1&0\\0&0\end{bmatrix},\qquad
  P_1=\begin{bmatrix}0&0\\0&1\end{bmatrix},\qquad
  S=\begin{bmatrix}0&1\\1&0\end{bmatrix}.
\]
The original rank-two field is constant, $H(k)=P_0$.  The pointwise unitary gauge is
$W(0)=I$ and $W(1)=S$.  Write $H^W(k)=W(k)^\dagger H(k)W(k)$, and let
$\mathcal T_0$ reconstruct a field after retaining only its trivial-character Fourier
coefficient.

\begin{theorem}[Two-point gauge--truncation noncommutation]
\label{thm:two-point-gauge-truncation}
For the exact field and gauge above, the original odd coefficient is zero, whereas
\[
  \widehat{H^W}_{\mathrm{odd}}
  =\frac12\begin{bmatrix}1&0\\0&-1\end{bmatrix}\neq 0.
\]
Moreover,
\[
  \mathcal T_0(H^W)\neq (\mathcal T_0 H)^W.
\]
\end{theorem}

\begin{proof}
Conjugation by $S$ exchanges $P_0$ and $P_1$.  Hence $H^W(0)=P_0$ and
$H^W(1)=P_1$.  The nontrivial character of $\mathbb Z/2\mathbb Z$ takes values $1$
and $-1$ at the two points, so the normalized odd coefficient is
$(P_0-P_1)/2$, while the same coefficient of the constant field vanishes.  Zero-mode
truncation fixes $H$.  If the two displayed routes were equal, evaluating their
zero-mode reconstruction at both momenta would force $P_0=P_1$, contradicting their
$(0,0)$ entries.  This is the checked proof sequence
\path{BG-FOURIER-THM-055}--\path{BG-FOURIER-THM-066}; the final declaration is
\path{ProjectKoios.BandGauge.Fourier.pointwiseSwapConjugation_zeroModeFourierTruncation_ne}.
\end{proof}

The theorem is a finite existence counterexample, not a universal statement about every
momentum-dependent gauge.  It establishes neither continuum locality nor decay,
smoothness, mesh convergence, runtime correctness, topology, or material adequacy.  In
rank one, scalar operator conjugation is trivial, so this rank-two witness must not be
reinterpreted as rank-one gauge-locality behavior.  The source module for the following
crosswalk is
\path{ProjectKoios/BandGauge/Fourier/GaugeTruncationCounterexample.lean}.

\begin{table}[ht]
\centering
\footnotesize
\begin{tabularx}{\textwidth}{@{}l >{\raggedright\arraybackslash}X@{}}
\toprule
ID & Exact Lean declaration and proof dossier \\
\midrule
\path{BG-FOURIER-THM-062}
  & \path{ProjectKoios.BandGauge.Fourier.pointwiseSwapConjugation_odd_coefficient}
    \newline \path{proofs/BandGauge/Fourier/BG-FOURIER-THM-062-gauged-odd-coefficient.tex} \\
\path{BG-FOURIER-THM-064}
  & \path{ProjectKoios.BandGauge.Fourier.pointwiseSwapConjugation_odd_coefficient_ne_zero}
    \newline \path{proofs/BandGauge/Fourier/BG-FOURIER-THM-064-gauged-odd-mode-nonzero.tex} \\
\path{BG-FOURIER-THM-066}
  & \path{ProjectKoios.BandGauge.Fourier.pointwiseSwapConjugation_zeroModeFourierTruncation_ne}
    \newline \path{proofs/BandGauge/Fourier/BG-FOURIER-THM-066-gauge-truncation-noncommutation.tex} \\
\bottomrule
\end{tabularx}
\caption{Key declarations for the exact counterexample.}
\label{tab:gauge-truncation-crosswalk}
\end{table}

\section{Exact finite truncation tails}
\label{sec:finite-truncation-tails}

Let $G$ be a finite additive group, let $\widehat G$ be its finite character group, and
let $c:\widehat G\to V$ be coefficients in a complex normed space.  For a declared
support $A\subseteq\widehat G$, write $c_A$ for the coefficients obtained by setting
all modes outside $A$ to zero, and define the deliberate-truncation residual
\[
  R_A(k)=\mathcal F^{-1}c(k)-\mathcal F^{-1}c_A(k).
\]
This is the interpolant-minus-truncated channel.  It is not the
reference-minus-interpolant channel and is not, by itself, the total residual.

\begin{theorem}[Finite Fourier-tail identities and bounds]
\label{thm:finite-fourier-tail-bounds}
Under the repository's finite Fourier normalization, the following hold.
\begin{enumerate}
  \item If $V$ is a complex Hilbert space, then
  \[
    \left(\frac1{|G|}\sum_{k\in G}\lVert R_A(k)\rVert^2\right)^{1/2}
    =\left(\sum_{\chi\notin A}\lVert c(\chi)\rVert^2\right)^{1/2}.
  \]
  \item For complex matrix coefficients, the same identity holds with the explicit
  entrywise Frobenius power
  $\sum_{i,j}|R_A(k)_{ij}|^2$ and
  $\sum_{i,j}|c(\chi)_{ij}|^2$.  The Lean root-mean-square theorems in items 1 and 2
  carry an explicit nonempty finite-domain premise.
  \item For any complex normed space and every $k\in G$,
  \[
    \lVert R_A(k)\rVert
    \leq \sum_{\chi\notin A}\lVert c(\chi)\rVert,
  \]
  and therefore
  \[
    \max_{k\in G}\lVert R_A(k)\rVert
    \leq \sum_{\chi\notin A}\lVert c(\chi)\rVert.
  \]
\end{enumerate}
\end{theorem}

\begin{proof}
The residual is definitionally the previously formalized finite Fourier truncation
residual.  Finite Parseval gives total residual power as $|G|$ times omitted coefficient
power; division by $\sqrt{|G|}$ yields the root-mean-square identities.  Pointwise, the
omitted-mode reconstruction formula and the triangle inequality bound the residual by
the sum of the norms of its summands.  Character phases disappear because each has norm
one.  The finite maximum statement is the coordinatewise supremum lift of this same
bound.  These steps are declarations \path{BG-FOURIER-THM-067}--
\path{BG-FOURIER-THM-075}.
\end{proof}

The matrix root-mean-square statement uses explicit entrywise Frobenius power, not the
ambient matrix norm.  Conversely, the maximum theorem uses the ambient norm on an
arbitrary complex normed space; no entrywise matrix-Frobenius maximum specialization is
claimed.  The maximum theorem requires no separate \path{Nonempty G} argument beyond
the additive-group structure, which already supplies $0:G$; this is not an empty-domain
claim.  None of the results binds floating-point coefficients, a runtime interpolant, a
mesh, physical hopping range, or a scientific threshold.  They also do not assert scalar
additivity of interpolation, truncation, and total errors.  The source module for both
following crosswalks is
\path{ProjectKoios/BandGauge/Fourier/ResidualChannels.lean}.

\begin{table}[ht]
\centering
\footnotesize
\begin{tabularx}{\textwidth}{@{}l >{\raggedright\arraybackslash}X@{}}
\toprule
ID & Exact Lean declaration and proof dossier \\
\midrule
\path{BG-FOURIER-THM-067}
  & \path{ProjectKoios.BandGauge.Fourier.truncationResidual_fourierReconstruction_eq}
    \newline \path{proofs/BandGauge/Fourier/BG-FOURIER-THM-067-truncation-channel.tex} \\
\path{BG-FOURIER-THM-068}
  & \path{ProjectKoios.BandGauge.Fourier.finiteFourierPower_truncationResidual_eq_omitted}
    \newline \path{proofs/BandGauge/Fourier/BG-FOURIER-THM-068-hilbert-tail.tex} \\
\path{BG-FOURIER-THM-069}
  & \path{ProjectKoios.BandGauge.Fourier.finiteMatrixFrobeniusPower_truncationResidual_eq_omitted}
    \newline \path{proofs/BandGauge/Fourier/BG-FOURIER-THM-069-matrix-tail.tex} \\
\path{BG-FOURIER-THM-070}
  & \path{ProjectKoios.BandGauge.Fourier.finiteRootMeanSquareNorm_truncationResidual_eq_sqrt_omitted}
    \newline \path{proofs/BandGauge/Fourier/BG-FOURIER-THM-070-hilbert-rms-tail.tex} \\
\path{BG-FOURIER-THM-071}
  & \path{ProjectKoios.BandGauge.Fourier.finiteRootMeanSquareNorm_truncationResidual_sq_eq_omitted}
    \newline \path{proofs/BandGauge/Fourier/BG-FOURIER-THM-071-hilbert-rms-squared-tail.tex} \\
\path{BG-FOURIER-THM-072}
  & \path{ProjectKoios.BandGauge.Fourier.finiteMatrixFrobeniusRootMeanSquare_truncationResidual_eq_sqrt_omitted}
    \newline \path{proofs/BandGauge/Fourier/BG-FOURIER-THM-072-matrix-frobenius-rms-tail.tex} \\
\path{BG-FOURIER-THM-073}
  & \path{ProjectKoios.BandGauge.Fourier.finiteMatrixFrobeniusRootMeanSquare_truncationResidual_sq_eq_omitted}
    \newline \path{proofs/BandGauge/Fourier/BG-FOURIER-THM-073-matrix-frobenius-rms-squared-tail.tex} \\
\bottomrule
\end{tabularx}
\caption{Exact power and root-mean-square declaration crosswalk.}
\label{tab:rms-tail-crosswalk}
\end{table}

\begin{table}[ht]
\centering
\footnotesize
\begin{tabularx}{\textwidth}{@{}l >{\raggedright\arraybackslash}X@{}}
\toprule
ID & Exact Lean declaration and proof dossier \\
\midrule
\path{BG-FOURIER-THM-074}
  & \path{ProjectKoios.BandGauge.Fourier.norm_truncationResidual_fourierReconstruction_le_omitted}
    \newline \path{proofs/BandGauge/Fourier/BG-FOURIER-THM-074-pointwise-maximum-tail.tex} \\
\path{BG-FOURIER-THM-075}
  & \path{ProjectKoios.BandGauge.Fourier.finiteMaximumNorm_truncationResidual_le_omitted}
    \newline \path{proofs/BandGauge/Fourier/BG-FOURIER-THM-075-finite-maximum-tail.tex} \\
\bottomrule
\end{tabularx}
\caption{Exact pointwise and finite-maximum declaration crosswalk.}
\label{tab:maximum-tail-crosswalk}
\end{table}
\FloatBarrier

\section{Finite route algebra}

The finite-lattice route layer defines coordinate pullback as $U^{\dagger}AU$.
Declarations \decl{SS-ROUTE-THM-001--010} prove compatibility with addition,
subtraction, composition, scalar-reference removal, Hermiticity, spin lift, folding, and
aligned defect extraction under their explicit isometry and identification premises.
They establish exact algebraic commutativity for the same declared maps. Nonlinear
fitting, separately selected gauges, and different model classes define different
routes and require separate analyses.
